\documentclass[10pt,a4paper]{amsart}
\usepackage[margin=19mm]{geometry}
\usepackage{amsmath,amssymb,mathtools}
\usepackage[scr=rsfs,cal=euler]{mathalfa}
\usepackage{xfrac}
\usepackage[unicode,hidelinks,bookmarks=false]{hyperref}
\newcommand{\ie}{i.e.\ }
\newcommand{\cf}{cf.\ }
\newcommand{\resp}{respectively\ }
\newcommand{\Subsection}{\subsection}
\providecommand{\nobreakdash}{\nobreak\hbox{-}}
\setlength{\emergencystretch}{2em}
\multlinegap0pt
\usepackage{mathtools}
\usepackage{amssymb}
\usepackage{xcolor}
\usepackage{enumerate}
\usepackage[shortlabels]{enumitem}
\usepackage{tikz}
\newtheorem{theorem}{Theorem}
\title{A hyperbolic model for two-layer thin film flow with a perfectly soluble anti-surfactant}
\author{Rahul Barthwal, Christian Rohde}
\date{Source: arXiv:2502.17205v2 (CC BY 4.0)}
\begin{document}
\maketitle

\noindent\emph{Source and licence.} A hyperbolic model for two-layer thin film flow with a perfectly soluble anti-surfactant. Version arXiv:2502.17205v2 (CC BY 4.0), distributed by arXiv under the Creative Commons Attribution 4.0 International licence, http://creativecommons.org/licenses/by/4.0/, as recorded in the arXiv licence metadata for this exact version and shown on the abstract page https://arxiv.org/abs/2502.17205v2. Full licence notice: \texttt{licenses/arxiv-2502.17205-CC-BY-4.0.txt}.\par\smallskip

\noindent\textit{Editorial scope.} Counted unit: the article's theorem theo1 (source0.tex lines 1077--1088). The article's complete original proof of it is delivered with it (source0.tex lines 1089--1110, one \texttt{proof} environment of 22 lines). No mathematics is written, repaired or re-derived: the statement and the proof are the article's own lines, in the article's own notation, and no retained line is substituted or re-flowed.  Retained uncounted as the source-local context the proof needs: lines 915--922; lines 1012--1015; lines 1064--1070.  Nothing was omitted from any retained span, so the package carries no omission record.  No retained line contains a citation, so the wrapper carries no bibliography and no bibliography entry.  No retained line contains a figure environment, an image-inclusion command or drawing code, so no image is delivered and the package declares no asset dependency.  Editorial formatting only, in addition: an AMS \texttt{amsart} preamble that restates the article's own declarations and macros used by the retained lines, namely the environments \texttt{theorem} on separate counters, together with the standard packages those lines need.  The retained environments print in this document as Theorem 1 in order of appearance, whereas the article prints the same items with the numbering of its own full document.  The complete native source of the article as distributed in the arXiv e-print for this exact version is bundled unchanged as \texttt{sources/173783/source0.tex}.
\begin{align}\label{eq: Main_system}
%\begin{cases}
    \dfrac{\partial f}{\partial t}+\dfrac{\partial}{\partial x}\left(\dfrac{1}{2}f^2b\right)&=0,\vspace{0.2 cm} \nonumber\\
    \dfrac{\partial b}{\partial t}+\dfrac{\partial}{\partial x}\left(\dfrac{1}{2}fb^2\right)&=0,\vspace{0.2 cm} \\
     \dfrac{\partial g}{\partial t}+\dfrac{\partial}{\partial x}\left(\dfrac{g^2q}{2}+fgb\right)&=0,\vspace{0.2 cm} \nonumber\\
  \dfrac{\partial q}{\partial t}+\dfrac{\partial}{\partial x}\left(\dfrac{gq^2}{2}+fbq\right)&=0\nonumber
 % \end{cases}
\end{align}

\begin{equation}\label{entropy}
\nabla_{\mathbf{U}} E(\mathbf{U})^\top\mathbf{DF(U)}=\nabla_{\mathbf{U}}Q(\mathbf{U})^\top
 \text{ for all $\mathbf{U}\in \mathcal U$}.
\end{equation}

\begin{equation}\label{Rtransformed}
\begin{array}{rclrcl}
   \xi_t+\dfrac{u}{2} \xi_x&=&0,&  u_t+\dfrac{3 u}{2} u_x&=&0, \\[1.5ex]
   \tau_t+\left(u+\dfrac{v}{2}\right)\tau_x&=&0, &\quad 
\eta_t+\left(u+\dfrac{3v}{2}\right)\eta_x&=&0.
\end{array}    
\end{equation}

\begin{theorem}[Entropy/entropy-flux pairs] \label{theo1}Let $\rho, \mu$ and $\nu$ be arbitrary functions in  $C^1(0,\infty)$.
For the system \eqref{eq: Main_system}, there exists a class of entropy/entropy-flux pairs $( E, Q)$ which is 
given by
\begin{equation}\label{eq:entropyclass}
\begin{array}{rcl}
 E[f, b, g, q]&=& \rho(fb)+\sqrt{fb}\mu\left(\dfrac{b}{\sqrt{fb}}\right)+\sqrt{gq}\nu\left(\dfrac{q}{\sqrt{gq}}\right)+\dfrac{1}{fb+gq},\\[3ex]
 Q[f, b, g, q]&=&\psi(fb)+(fb)^{3/2} \mu\left(\dfrac{b}{\sqrt{fb}}\right)+(\sqrt{gq}) \nu\left(\dfrac{q}{\sqrt{gq}}\right)(fb+gq)\\
&&\hspace{3.5 cm}-\ln \left((fb+gq)^{3/2}\right)-\dfrac{fb}{2(fb+gq)}.
\end{array}
\end{equation}
The function $\psi$ is determined from $\psi'(w)=\dfrac{3}{2}  w\rho'(w)$.% with $u=fb$.
\end{theorem}

\begin{proof} We have to prove the compatibility relation \eqref{entropy} for $(E,Q)$ from \eqref{eq:entropyclass}.  In view of  the bijectivity of the map $(\xi, u, \tau, \eta)\mapsto (f, b, g, q)$, it suffices to  show that 
the transformed compatibility relation holds for the functions $E[\xi, u, \tau, \eta] := E[f, b, g, q]$ and
$Q[\xi, u, \tau, \eta] := Q[f, b, g, q]$.  From \eqref{Rtransformed} we observe that the componentwise 
relations are given by 
\begin{equation}\label{eq: 4.6Relations}
%\begin{cases}
Q_{\xi}=\dfrac{u}{2} E_{\xi},  \,  \,% \vspace{0.2 cm}\\
    Q_u=\dfrac{3u}{2}E_u,   \,\, %\vspace{0.2 cm}\\
    Q_{\tau}=\left(u+\dfrac{v}{2}\right) E_{\tau},    \, \, % \vspace{0.2 cm}\\\
    Q_{\eta}= \left(u+\dfrac{3v}{2}\right) E_{\eta}.
%\end{cases}    
\end{equation} 
The quantity $v= gq$ is related to the functions $u$ and $\eta$ by the nonlinear relation $u+v=\eta v^{1/4}$.

The functions $E$ and $Q$  from \eqref{eq:entropyclass}  in terms of the Riemann invariants $\xi, u, \tau, \eta$   are given by 
\[\begin{array}{rcl}
E[\xi, u, \tau, \eta]&=& \rho(u)+\sqrt{u}\mu(\xi)+\sqrt{v}\nu(\tau)+\dfrac{1}{\eta v^{1/4}},\\[2ex]
Q[\xi, u, \tau, \eta]&= &\psi(u)+\dfrac{u^{3/2}}{2}\mu(\xi)+\dfrac{\sqrt{v}\nu(\tau)}{2}\left( \eta v^{1/4}+u\right)-\dfrac{3}{2}\ln\left(\eta v^{1/4}\right)-\dfrac{u}{2\eta v^{1/4}}.
\end{array}
\]
We readily  check that the pair $\left(E[\xi, u, \tau, \eta], Q[\xi, u, \tau, \eta]\right)$ satisfies \eqref{eq: 4.6Relations}.
 \end{proof}

\end{document}
